\documentclass[11pt,a4paper]{article}
\usepackage[utf8]{inputenc}
\usepackage[margin=1in]{geometry}
\usepackage{xcolor}
\usepackage{hyperref}

\hypersetup{colorlinks=true, linkcolor=blue!70!black, urlcolor=blue!70!black}

\begin{document}

\noindent \textbf{Date:} October 29, 2026\\[1.5ex]
\textbf{To:}\\
Prof. Nicholas A. Palomares, Ph.D.\\
Editor-in-Chief, \textit{Journal of Language and Social Psychology}\\
SAGE Publications

\vspace{2ex}
\noindent \textbf{Subject:} Submission of Original Research Manuscript: \textit{``The Hypocrisy of Parity: Discursive Patterns of Victimhood Signaling and Self-Serving Attribution in Exploitative Relationships''}

\vspace{2ex}
\noindent Dear Prof. Palomares and the Editorial Board,

\vspace{1ex}
\noindent I am pleased to submit my original research manuscript entitled \textbf{``The Hypocrisy of Parity: Discursive Patterns of Victimhood Signaling and Self-Serving Attribution in Exploitative Relationships''} for consideration for publication in the \textit{Journal of Language and Social Psychology}.

\vspace{1.5ex}
\noindent \textbf{Significance and Disciplinary Fit with the \textit{Journal of Language and Social Psychology}:}\\
At the heart of discursive psychology, interpersonal communication, and language studies lies a central inquiry: how do individuals deploy language and rhetorical strategies to construct moral accountability, deflect culpability, and manage interpersonal relationships? In contemporary interpersonal relationships—spanning romantic partnerships, domestic cohabitation, family ties, and peer bonds—individuals increasingly adopt egalitarian discourse, frequently demanding strict 50/50 parity, emotional labor sharing, and structural fairness in the treatment of others. Paradoxically, an acute behavioral breakdown frequently unfolds in the actual treatment of others: individuals who most vocally champion parity and broadcast personal victimhood simultaneously engage in systematic, self-serving exploitation of those around them.

While traditional laboratory experiments and self-report surveys suffer from severe social desirability bias, our study introduces a computational mixed-methods architecture that leverages naturalistic online relationship conflict disclosures ($N = 142,580$ real-world dispute narratives from Reddit adjudication archives) as an unobtrusive observational window into authentic interpersonal language and behavior. This research directly bridges social psychological constructs—specifically \textbf{Moral Hypocrisy}, \textbf{Self-Serving Attribution}, and the \textbf{Dark Ego Vehicle Principle}—with empirical \textbf{Discursive Patterns} and account-making strategies, directly serving the core aims and scope of the \textit{Journal of Language and Social Psychology} and relevant discourse-analytic communities (such as \textit{Discourse Studies}).

\vspace{1.5ex}
\noindent \textbf{Methodological Rigor and Key Contributions:}
\begin{itemize}
    \item \textbf{Unobtrusive Naturalistic Discourse Corpus:} We analyze $N = 142,580$ authentic interpersonal conflict narratives, capturing real-world conversational accounts, rhetorical framing, and behavioral treatment of others rather than artificial survey responses.
    \item \textbf{Discursive Psychology \& Theoretical Integration:} Grounded in Discursive Psychology, Moral Hypocrisy Theory, Self-Serving Attribution Theory, and the Dark Ego Vehicle Principle, we specify and empirically validate a structural equation model (PLS-SEM) demonstrating that Egalitarian Vocabulary \& Parity Demands (\textit{EVP}) drive Exploitative Interpersonal Behaviors (\textit{EIR}) through Self-Serving Attribution \& Licensing (\textit{SSA}; indirect effect $\beta = 0.190, p < 0.001$), with Virtuous Victimhood Signaling (\textit{VVS}) serving as an offensive moderating lever ($\beta = 0.248, p < 0.001$).
    \item \textbf{Generative AI Cultural Informants for Pragmatic Decoding:} Deploying a triangulated multi-model Large Language Model pipeline (GPT-4o, Claude 3.5 Sonnet, Gemini 1.5 Pro) with strict temperature control ($T=0.10$), we demonstrate that LLMs possess superior pragmatic sensitivity in decoding latent discursive hypocrisy ($F1 = 0.893$) compared to traditional sentiment dictionaries ($F1 \le 0.612$).
    \item \textbf{Actionable Interpersonal Communication Insights:} We formulate concrete insights for interpersonal communication scholars and relationship practitioners to deconstruct how progressive psychological terminology is co-opted for interpersonal manipulation and moral licensing.
\end{itemize}

\vspace{1.5ex}
\noindent \textbf{Declarations and Ethics Compliance:}
\begin{enumerate}
    \item This manuscript is original, has not been published previously, and is not currently under consideration for publication elsewhere.
    \item The author has approved the manuscript and agrees with its submission to the \textit{Journal of Language and Social Psychology}.
    \item Data collection adhered strictly to the Association of Internet Researchers (AoIR) ethical guidelines; all user identifiers were cryptographically hashed (SHA-256) to ensure complete non-reidentifiability under IRB public exemption standards.
    \item AI assistance utilized during drafting (literature structuring and experimental LLM informant elicitation) is fully transparent and disclosed in accordance with ICMJE and COPE guidelines.
\end{enumerate}

\vspace{2ex}
\noindent Thank you very much for your time, consideration, and editorial stewardship. I look forward to receiving the reviewers' constructive evaluations.

\vspace{2.5ex}
\noindent Sincerely yours,\\[2ex]
\textbf{Peerapat Asoktummarungsr}\\
Corresponding Author\\
Email: \texttt{nuboat@gmail.com}

\end{document}
